\bookchapter{One source, one energy convention: isotope tests and the physical map}
\label{ch:common-map}

\subsection{The question after a complete construction}

The completed constructions now supply more than a roster of isotope labels. They retain legal internal states, occupied electronic sectors, field squares, counterterms, recoil, and the free systems against which formation is measured. This makes a more specific physical question possible: can one source law calculate the bound object and its free constituents in the same energy convention, and then carry that law across isotopes? The question connects the exact matter construction to the measured energy comparison. Its operands are the full systems developed in the preceding chapters.

Three recent records make the distinction concrete. The carbon isotope calculation transfers an already specified source into a complete prepared-electron family. The light-isotope calculation tests a particular scalar species-current law against a held helium extension. The 37-isotope preparation lays out a common prediction convention and preserves the structural operands for the next source derivation. These records have different completed outputs. Their relationship gives the volume a precise source-to-energy map.\samref{map-carbon}\samref{map-light}\samref{map-plan}

The organization in this chapter is an editorial synthesis of those completed records. The source hypotheses and numerical outcomes retain their original attribution. Sean Brady is the originator and conceptual director; OpenAI ChatGPT and Codex are AI research collaborators and co-authors. In particular, the scalar species-current extension is recorded as a Codex contribution, rather than attributed to the original matter ledger.

\subsection{The calculated atomic reference}

For a neutral isotope with proton count $Z$ and neutron count $N$, the prepared common-law comparison is
\begin{equation}
B_{\rm at}(Z,N)=Z E_{\mathrm{H1,free}}+N E_{\mathrm{n,free}}-E_{\rm atom}(Z,N).
\label{eq:atomic-binding-common}
\end{equation}
Each energy is calculated with the same source, electron, field, counterterm and energy-zero convention. A measured composite mass is a later comparator. The free endpoints are mathematical recipients in the same construction as the atom. This is particularly useful when a source contains a nonzero self-energy or an identity baseline: the term has a definite place in both sides of the subtraction.\samref{map-plan}

Common additive proton, neutron and electron rest baselines cancel once in Equation~\eqref{eq:atomic-binding-common}. The masses inside kinetic and collective recoil operators still affect the calculated state. Their role is dynamical. An independently supplied inertial mass must therefore remain listed among the inputs even when a common rest-energy constant disappears from the formation difference.

The neutral-atom convention also differs from bare-nuclear binding by a specified electronic contribution:
\begin{equation}
B_{\rm at}=B_{\rm nuclear}+B_{\rm electrons}(\mathrm{atom})-ZB_{\rm electrons}(\mathrm{H1}).
\end{equation}
Here the electronic binding quantities are positive. A comparison to a bare-nucleus reference needs this conversion and the corresponding bounds. It cannot change the reference system merely by changing the label on a table column. Fragment separation energies likewise come from differences of calculated endpoints. The 37-isotope preparation explicitly excludes measured light-fragment bindings as fixed offsets in the prediction inputs.

The same principle applies to numerical enclosures. Subtracting two variational upper witnesses does not establish an enclosure for their difference. A binding interval needs compatible lower and upper information for the complete systems. The source preparations retain whole-ground lower and upper recipients, as well as the reference-dependent electronic and recoil corrections needed when the nuclear source changes.

\subsection{Units and source identification are separate steps}

The retained independent unit map is
\begin{align}
E_{\rm unit}&=2\alpha^2m_ec^2,\\
L_{\rm unit}&=\frac{\hbar c}{2\alpha m_ec^2},\\
T_{\rm unit}&=\frac{\hbar}{E_{\rm unit}}.
\end{align}
The declared inputs are the existing CODATA reference values for the fine structure constant, electron rest energy, Planck constant, speed of light and elementary charge. These provide a common dimensional chart without selecting parameters from isotope binding values. The provenance of those constants remains part of the calculation.\samref{map-plan}

A dimensional chart and a source assignment do different work. The chart says how a quantity with an established role is expressed in energy, length or time units. A source assignment says which internal current acts on which field, with what spatial profile, radius and coupling. The retained D/C source and the electronic source have their own specified operators. Placing both beside the same unit ruler does not determine their intertwiner. That map is an explicit part of the continuing connector-action question.

This separation explains why a dimensionless spectrum remains valuable before all physical inputs have been selected. It can determine sectors, degeneracies, conditional minima, operator identities and source-dependent thresholds. A physical prediction then uses the selected source law and the common unit map. The complete earlier spectrum is retained as an operand, rather than recalculated under a new interpretation of its units.

\subsection{Carbon: a completed source-preserving isotope transfer}

The completed carbon calculation receives the neutral C4-rotor source with fixed $U=(12,0)$, $q=144$, $\gamma=\xi=1$, intrinsic coefficient $I=282$, and contact lowering $-12$. The complete source coefficients are
\begin{equation}
C_{0,12}=\frac{1536}{5},\qquad C_{0,13}=\frac{1476}{5}.
\end{equation}
The mass recipient has the form $\mu=B+C_0/r$, with the C13 inventory obtained by adding the same complete free-neutron reference to C12. Both existing C12 baseline inventories are retained. The source model supplies the isotope transfer without using a measured C13 mass to select its parameters.\samref{map-carbon}

The receiving electronic construction retains the common basis, complete massive kinetic operator, finite-core scalar and axial potential, full direct and exchange terms, collective momentum and the outside-span contribution. Both 20-spin and 24-spin banks are kept, with every electron count from zero through six. The declared homothetic control radius is $1/65536$; the independent physical-radius member and the whole source/field phase minimization remain separately identified.

The calculation is complete across 4,791 retained dilations and 9,582 new profiles, producing 134,148 envelopes. The source reports 670,740 independent exact checks. At the neutral 24-spin optimum both baselines select the unique sampled dilation $511/512$ and occupation mask $15728643$, which doubly occupies spatial indices 0, 10 and 11. The matching C12 static lower endpoints stay fixed. The C13 transfer reduces the recoil upper allowance.

\begin{table*}[htbp]
\centering\small
\caption{Prepared carbon attachment lower certificates in the retained nominal source chart.}
\begin{tabular}{lrrr}\toprule
Baseline & C12 (eV) & C13 (eV) & Gain (eV)\\\midrule
B1 inventory control &876.015801973&876.181678261&0.165876288\\
Reference free inventory &876.351683954&876.353893038&0.002209083777\\\bottomrule
\end{tabular}
\end{table*}

These values are prepared attachment certificates. Their change has an identified mechanism: the changed complete source mass tightens the recoil allowance. The physical isotope shift is a different observable. All best-grid dilations for the occupied charge families are unchanged. Five B1 point/charge pairs change the selected 24-spin spatial mask; the reference-inventory upper witness has no corresponding spatial change. The enumerator preserves every exact dilation tie while retaining one endpoint occupation witness at each point.

The 20-spin neutral control selects $841/1024$ and gives C13 certificates $458.214795799$ and $458.299935741$ eV in the two baselines. The comparison exposes two distinct improvements: enlarging the orbital family changes the attachment certificate substantially; transferring the isotope mass changes a smaller recoil contribution within that family. The source classification is: \textbf{The test result suggests strong contact with the concept.} Its tested concept is reuse of the original electron operands with the changed complete source mass and fresh all-charge occupation enumeration in the declared model.

The recorded scientific grid is complete. The original guest cold-closure receipt remains partial, and the supplemental local recovery restores all 178 stores, all 4,788 cloud point references, 9,932 mathematical values and 44,642 objects, including 178 actual worker-summary reads with recomputation disabled. Those receipts explain the completed result without converting the earlier partial closure into an altered historical record.\samref{map-carbon}

\subsection{A legal-state neutron current changes the existing source}

The separate neutron-to-electron test starts with the original six-C4 intrinsic operator. Its conditional minima at magnetizations $-6,-4,-2,0,2,4,6$ are
\begin{equation}
626,\quad110,\quad110,\quad110,\quad110,\quad110,\quad626.
\end{equation}
This allows an exact reduction of the complete ground question to 49 H2 tuples and 343 tuples each for H3 and He3, retaining the two independent Gray halves. The source report records 8,820 Bernstein coefficients agreeing with the literal formulas: 8,748 positive, 72 zero and none negative in the tested parameter box.\samref{map-light}

For the scalar-current hypothesis, the actual branches are H2$(4,4)$, H3$(2,4,4)$ and He3$(4,4,4)$, with their global-sign partners. H3 can lower its proton current to two while retaining its intrinsic minimum. Adding a neutron can therefore change a current already present in the source. A fixed current increment for the new constituent would omit the relaxation of the existing constituent.

The retained shifted-profile electronic potential is
\begin{align}
V(x)&=-\sum_{p}G_R(x-x_v)
-\frac1{10}\sum_v a_v\,U_v\!\cdot n(x)\,G_R(x-x_v),\\
G_R(y)&=\frac1{\max(|y|,R)},\qquad a_p=1,\quad a_n=\eta.
\end{align}
Each shifted profile remains in the calculation. The aligned branch has total-current coefficients $4+4\eta$, $2+8\eta$ and $8+4\eta$ for H2, H3 and He3 respectively. These coefficients summarize the co-centred or far-field leading term; they do not erase the finite spatial source.

The light-source diagnostic calibrates the scalar family to the retained nominal H2/H3/He3 atomic formation values. It finds the local enclosures
\begin{align}
\eta&\in[0.6479399,0.6479401],\\
\gamma&\in[0.337148885323,0.337149328220],\\
S_{\rm MeV}&\in[11.908003512263,11.908055649761].
\end{align}
These are source-only calibration enclosures. They have measured binding contact and remain outside the inputs of the prepared no-binding-calibration prediction. Fresh atomic and recoil intervals are carried separately. The calculation completed 25 native calls with no scientific-call failure. The earlier all-current-four branch calculation remains labeled as preliminary and is superseded by the legal-state branches above.

\subsection{The held helium extension identifies the scalar law's outcome}

The held extension uses the same intrinsic/free/current law on the original He4 outer K4 inventory. The tested geometry is an explicitly assigned regular tetrahedral tangent shell with one PP, four PN and one NN contact. Four legal templates supply a complete source-energy upper witness and the binding coefficient lower
\begin{equation}
\begin{aligned}
g_4&\geq16\bigl[(\gamma-0.1)(1+4\eta+\eta^2)\\
&\qquad-0.1(2+2\eta^2)\bigr]-1\\
&\geq9.589406784.
\end{aligned}
\end{equation}
The actual H2 branch has $g_2\leq0.2078544$. Combining its independent atomic reference, the new atomic correction and the same source scale gives
\begin{equation}
B_{\rm at}(\mathrm{He4})\geq102.4621574335\ \mathrm{MeV}.
\end{equation}
The retained reference is $28.295662378$ MeV. The lower bound lies $74.1664950555$ MeV above that reference. The complete electron/recoil allowances are included. For the specified static scalar-current branch box and equal-shell complete-K4 extension, the recorded classification is: \textbf{The test falsifies the concept.}\samref{map-light}

The outcome is obtained from a compatible full-energy witness; a further minimization would increase the binding lower bound. This is why the test could settle its stated scalar family without another large ground-state enumeration. The exact legal-state receiver and the observed current-two branch remain completed structural results. The retained C13 electron expansion has its own fixed source and remains unchanged by this separate scalar-law test.

\subsection{A prepared 37-isotope recipient, with an explicit next operand}

The staged roster contains 31 primary isotopes in six element chains: O16--18; Ca40,42,43,44,46; Fe54,56,57,58; Ni58,60,61,62,64; Sn112,114,115,116,117,118,119,120,122,124; and Pb204,206,207,208. C12 and C13 form the carbon pair. H2, H3, He3 and He4 form the separate light diagnostic panel. H1 and the free neutron supply the calculated references. These groups total 37 prediction targets; the larger tensor identity roster is retained for subsequent expansion under the same law.\samref{map-plan}

The preparation is complete; no new physical prediction or native scientific calculation is recorded for that panel. Its immediate operand is the connector action and the common source/current/radius map. The retained PN actions are $38/36$, the PP actions $10/8$, with roles $60/0/1$ and the actual H3 neutron-neutron connector. Both recovered PN and PP connectors have $J=1$. That identity is an operator property; selecting the physical coefficients still requires the same action and source-current map.

The He3 fragment recipient makes the next question explicit:
\begin{equation}
6\tau_h-\epsilon j+P_{\rm paid}+\Delta_{\rm changed}<0.
\end{equation}
Here the hopping clock, source coefficient, paid preparation and changed complete reference each retain their source-defined meanings. A positive-gate construction carries its mandatory identity baseline once. The original two polarizations, field squares, self and mixed counterterms, and the full free reference remain in that expression. The 48 placement, 192 fragment, 144 shared-cell control, 64 colour and complete 16-state D/C spaces are retained as differently typed operands.\samref{map-plan}\samref{NUC_FRAGMENT}

The carbon source trial supplies an earlier conditional prediction precedent. Using independent C12 radius data and declared unit-contact choices, it froze a whole-ground separation interval $[6.983342570476823,7.007083582298198]$ MeV before comparison to $4.9463087(5)$ MeV. The prepared panel preserves that result and its information order. It excludes isotope-specific offsets, binding-selected radii and the fitted scalar parameters above from the new prediction inputs.

The matter-side achievement is thus a continuous set of recipients: legal source states, complete energy differences, electronic occupation and recoil, declared physical units, and a common-law isotope comparison. The remaining connector/current/radius assignment is stated as an exact source question. It can be developed from the retained operands while every completed structural, atomic and comparative result remains available to the book.
